% Copyright 2026 The terCAD team. All rights reserved.
% Use of this content is governed by a CC BY-NC-ND 4.0 license that can be found in the LICENSE file.

\subsection{Financial Statements}
\markboth{Budgeting}{Financial Statements}

In accounting, two reports answer most of the questions about the state of finances: the \emph{Balance Sheet} and the \emph{Income Statement}. Both are built from the same source -- the full log of recorded operations, known as the \emph{General Ledger} -- but they look at it from different angles \cite{Hart22}.

\begin{description}
  \item[Balance Sheet] is a snapshot at a given moment: what is owned (\emph{assets}), what is owed (\emph{liabilities}), and the difference between them (\emph{equity}, or net worth). It follows the accounting equation:
  \[
    \textit{Assets} = \textit{Liabilities} + \textit{Equity}
  \]
  \item[Income Statement] covers a period of time: all incomes received minus all expenses made, which results in a net profit (savings) or a net loss (overspending):
  \[
    \textit{Net Result} = \sum \textit{Incomes} - \sum \textit{Expenses}
  \]
\end{description}

\noindent Why bother with them when it comes to personal finances? Because a quick look at
these figures shows whether things are moving in the right direction. Drilling down into the
details helps you to identify the categories or months that need more attention.

\subsubsection{Chart of Accounts}

Businesses organize the general ledger using a chart of accounts. This is a list of all the
accounts and categories in use, usually with numeric codes that determine their position on
reports (e.g. assets, liabilities and equity on the balance sheet; revenues and expenses on
the income statement).

Fingrom retains this concept, but without codes. Its chart of accounts comprises two lists
that will be familiar from the previous section:

\begin{center}
\begin{tabularx}{\textwidth}{|l|X|l|}
  \hline
  \textbf{Accounting} & \textbf{Fingrom} & \textbf{Report} \\
  \hline
  Assets & Accounts: cash, bank accounts, debit cards, deposits, investments & Balance Sheet \\
  \hline
  Liabilities & Accounts: credit cards, loans, mortgages (negative balance) & Balance Sheet \\
  \hline
  Revenues & Incomes, attributed to the account that received them & Income Statement \\
  \hline
  Expenses & Bills, attributed to budget categories & Income Statement \\
  \hline
\end{tabularx}
\end{center}


\subsubsection{Cash and Accrual Methods}

There are two methods of recording financial transactions. The accrual method registers
income when it is earned and expenses when they are incurred, regardless of when the money
actually changes hands. It requires additional accounts such as receivables, payables, and
prepaid or accrued expenses. The cash method only registers an operation when money changes
hands.

Fingrom uses the cash method, whereby a bill or income is recorded when it is paid or
received. This means that the balance of each account always reflects the actual amount of
money held in it. This makes reports more streamlined and easier to verify against a bank
statement. Rather than being recorded, upcoming obligations are planned: recurring payments
(\cref{img:recurring-payments}) create bills automatically when they are due, and irregular
expenses can be smoothed with yearly categories or monthly coefficients
(see \cref{budget:coef}).

\img{img/ch01-s02-automation__recurring-payments}{Recurring Payments}{img:recurring-payments}

\noindent In this sense, the list of accounts on the main page, along with the total,
constitutes the application's balance sheet, while the \textbf{Summary} table described below
constitutes its income statement.


\subsection{Summary Table}
\label{summary:table}
\markboth{Budgeting}{Summary Table}

On a wide screen, the main page of Fingrom shows an overview of the last six months: how much
was spent in each budget category, how much was received on each account, and what remained
at the end of each month (see \cref{img:summary-table}).

\img{img/ch01-s02-main__summary}{Summary Table}{img:summary-table}


\subsubsection{When It Is Shown}

The default layout of the main page depends on the size of the window. The
wide-screen layout is used when the window is at least 1440 pixels wide and 480 pixels high
-- a size typical of desktop monitors, laptops and tablets in landscape mode. Narrower
windows display lists of accounts, budgets and bills instead (see \cref{img:main-page}).\\
\\

\noindent The wide-screen layout is a grid of 12 columns by 7 rows:

\begin{center}
\begin{tabular}{|p{0.27\textwidth}|p{0.27\textwidth}|p{0.27\textwidth}|}
  \hline
  \multicolumn{3}{|c|}{} \\
  \multicolumn{3}{|c|}{\textbf{Summary} (rows 1--4, full width)} \\
  \multicolumn{3}{|c|}{} \\
  \hline
  \centering \textbf{Candlestick (OHLC) Chart} &
  \centering \textbf{YTD Expenses} &
  \centering \textbf{Forecast Chart} \tabularnewline
  \centering accounts & \centering bills & \centering budgets \tabularnewline
  \hline
\end{tabular}
\end{center}

\noindent The upper part (four rows of the grid) is occupied by the \textbf{Summary} table,
and the lower part (three rows) is divided equally between three charts: the
\textbf{Candlestick (OHLC) Chart} of accounts, the \textbf{YTD Expenses} of bills, and the
\textbf{Forecast Chart} of budgets (see \cref{budget:metrics}).

This layout is just the default one. The \textbf{Customize Page} action
(\cref{img:main-editor}) on the main page allows you to rearrange, resize, add or remove
components, and the configuration is stored separately for each window size. This means that
even the \textbf{Summary} table can be shown on a mobile device. Once a custom configuration
has been saved, it replaces the default one.

\img{img/ch01-s02-main__editor}{Customize Main Page}{img:main-editor}


\subsubsection{Structure of the Table}

The table has eight columns: the title; one column for each of the last six months, with the
current month being the rightmost; and the Summary column, which shows the total for these
six months. The rows are organized into three blocks:

\begin{description}
  \item[Budgets] -- one row per active budget category with the amount spent in it in each
  month. This block is the expense part of the Income Statement.
  \item[Accounts] -- one row per active account that received at least one income, with the
  incomes per month. Transfers between own accounts are not incomes and are therefore not
  counted. This block is the revenue part of the Income Statement.
  \item[Total] -- the net result of each month: all incomes minus all bills. A positive value
  means that money was saved; a negative one means that more was spent than earned.
\end{description}

\noindent A simplified example (with only three months shown):

\begin{center}
\footnotesize
\setlength{\tabcolsep}{4pt}
\begin{tabular}{|l|r|r|r|r|}
  \hline
   & \textbf{Aug} & \textbf{Sep} & \textbf{Oct} & \textbf{Summary} \\
  \hline
  \textbf{Budgets} & & & & \\
  - Food/Groceries & 410.00 & 385.50 & 220.10 & 1\,015.60 \\
  - Housing/Rent & 900.00 & 900.00 & 900.00 & 2\,700.00 \\
  - Leisure/Hobbies & 120.00 & 260.00 & 40.00 & 420.00 \\
  \textbf{Accounts} & & & & \\
  - Bank/Salary & 2\,000.00 & 2\,000.00 & 2\,000.00 & 6\,000.00 \\
  - Cash & 0.00 & 150.00 & 0.00 & 150.00 \\
  \hline
  \textbf{Total} & \textbf{570.00} & \textbf{604.50} & \textbf{839.90} & \textbf{2\,014.40} \\
  \hline
\end{tabular}
\end{center}

\subsubsection{How the Values Are Calculated}

\textbf{Months follow the start of month setting.} For example, if the start of month is set
to 25, the October column covers the period from 25 October to 24 November. The current
period is always the rightmost column.

\textbf{All amounts are in the default currency.} Bills and incomes in different currencies
are converted using exchange rates so that the figures can be compared and totalled directly.

\textbf{Deactivated items are hidden.} Deactivated budget categories and accounts are not
given their own rows.

\textbf{Order follows the sorting preference.} The order in which categories and accounts are
listed is the same as that selected for lists elsewhere in the application (by totals, dates
or titles).

\textbf{The table is always up to date.} It is recalculated with every change, meaning a new
bill appears in the current column immediately.

\noindent Please note that the 'Total' row is calculated from all bills and incomes of the
month, including those from deactivated categories and categories marked 'Exclude from
Totals'. Therefore, the total may differ from a manual sum of the visible rows, but it always
reflects the actual cash flow.


\subsubsection{Reading the Table}

The table is the starting point for the monthly review recommended in the previous section:

\begin{enumerate}
  \item \textbf{Check the Total row first.} Even if every single category seems reasonable, a
  sequence of negative months is the clearest sign that the budget does not match the income.
  \item \textbf{Look for trends along each row.} If the growth of a category is justified, it
  needs either a higher limit or more attention.
  \item \textbf{Compare with the limits.} The Summary column divided by six gives the average
  monthly spending. This provides a natural basis for incremental budgeting and adjusting the
  Limit per Month.
  \item \textbf{Check the stability of incomes.} Inconsistent amounts in the 'Accounts' block
  suggest that some categories should be assigned relative limits so that they adapt
  automatically to good and bad months.
  \item \textbf{Keep in mind the current month.} The values in the rightmost column are
  naturally lower as it is still in progress. The Forecast Chart below the table shows where
  the current month is heading.
\end{enumerate}
